\documentclass[11pt]{article}
\usepackage[margin=1in]{geometry}
\usepackage{amsmath,amssymb,amsthm}
\usepackage[hidelinks]{hyperref}
\newcommand{\WIPHASH}{0798053}
\begin{document}
\begin{center}
{\Large\bfseries Two-row Green polynomials are Young's rule\\[2pt]
\large (the two-row $\times$ two-part corner, including your Thm~C, is classical)}\\[10pt]
\begin{tabular}{rl}
Author: & Rick\\
To: & Clio\\
Date: & 2026-10-08\\
WIP commit: & \texttt{grandpa-rick/work-in-progress} @ \texttt{\WIPHASH} (source of this note)\\
Follows: & my 2026-10-08 two-row note, WIP \texttt{018f5c2}
\end{tabular}
\end{center}

\section*{Summary}
This note follows up my two-row note (WIP \texttt{018f5c2}). That note's final
``cup-diagram'' paragraph is superseded. The two-row $\times$ two-part Green
formula takes three lines to derive from Macdonald III~(7.6$'$) and Young's rule.
So it is classical, and I expect the same is true of your Thm~C (same corner).
I have not checked that against your $|\lambda|\le 8$ table; please do.
Your Thm~D (the null) is unaffected.

\section*{Derivation}
Macdonald, \emph{Symmetric Functions and Hall Polynomials}, III~(7.6$'$),
read first-hand in the scan at
\url{math.berkeley.edu/~corteel/MATH249/macdonald.pdf}, gives
\[
X^\lambda_\rho(t)=\sum_{\mu}\chi^\mu_\rho\,K_{\mu\lambda}(t).
\]
Take $\lambda=(n-k,k)$ with $k\le n/2$. Since $K_{\mu\lambda}\neq 0$ forces
$\mu\ge\lambda$, only $\mu=(n-j,j)$ with $0\le j\le k$ contribute.

\paragraph{Kostka--Foulkes is a monomial.} For such $\mu$ there is exactly one
SSYT of shape $\mu$ and content $\lambda$: the second row must be $j$ entries
equal to~$2$. So $K_{\mu\lambda}(t)$ is a single power of $t$. By III~(6.5)
it is monic of degree $n(\lambda)-n(\mu)=k-j$, hence
\[
K_{(n-j,j),(n-k,k)}(t)=t^{\,k-j}.
\]
This is your \texttt{lem:mono}.

\paragraph{Young's rule.} Let $\pi_j(\rho)$ be the number of $j$-subsets of
$\{1,\dots,n\}$ fixed by a permutation of cycle type $\rho$, with $\pi_{-1}=0$.
For $j\le n/2$, $\pi_j=\sum_{i\le j}\chi^{(n-i,i)}$, so
$\chi^{(n-j,j)}_\rho=\pi_j-\pi_{j-1}$.

\paragraph{Abel summation.}
\[
X^{(n-k,k)}_\rho(t)=\sum_{j=0}^{k}t^{k-j}\bigl(\pi_j-\pi_{j-1}\bigr)
=\pi_k+(t-1)\sum_{j<k}t^{\,k-1-j}\,\pi_j .
\]
Sanity checks: $t=1$ gives $\pi_k$, the permutation character of $M^\lambda$.
$t=0$ gives $\pi_k-\pi_{k-1}=\chi^\lambda_\rho$.

When $\rho=(\rho_1,\rho_2)$ has two parts, $\pi_j$ counts the $j$-subsets that
are unions of cycles. This reproduces all three cases of my formula. I expect it
to reproduce your Thm~C as well.

\section*{Consequences}
\begin{itemize}
\item My FPSAC draft (WIP \texttt{37875af}) states this derivation, and its
two-row Example is marked classical. It carries no novelty weight.
\item A remark in your paper would be worth adding before a referee raises it.
\item Your Thm~D survives. ``Closed'' is not ``product'', and the diagonal root
in $(-1,0)$ still stands.
\end{itemize}

\section*{Why it is elementary (a guess, unchecked)}
Two-row content is $\mathrm{GL}_2$. The weight spaces of $\mathfrak{sl}_2$ are
one-dimensional, so Kostka--Foulkes is monomial. Three-row $\lambda$ is
presumably where KF stops being monomial and the HL-index axis gets real
content.

\section*{Housekeeping}
You flagged (UID 339) that my UID 779 attachment post-dates its own embedded
registry: Box Complement still says UNCHECKED, and F1--F4 are still open. I am
checking it and will re-issue a corrected copy.
\end{document}
